\section{Vorstellung Kronos}

\begin{frame}{Kronos: Kinematic Robotic Operations \& Navigation Orchestration System}
    \begin{alertblock}{Überblick}
        Kronos ist die Kinematik- und Pfadplanungs-Engine unseres 6-Achs-Slicers. 
        Es übersetzt die Druckpfade von Medusa in präzise Bewegungsbefehle für den Roboterarm, 
        um eine genaue und effiziente Fertigung zu gewährleisten.
    \end{alertblock}

    \vspace{0.25cm}
    \textbf{Funktionsweise von Kronos:}
    \begin{itemize}
        \item \textbf{Inverse Kinematik (IK):} Berechnet die 6 Gelenkwinkel $(J1...J6)$ für eine gegebene TCP-Position und Orientierung.
        \item \textbf{Pfad-Interpolation:} Stellt sicher, dass Bewegungen zwischen zwei Punkten linear im kartesischen Raum sind.
        \item \textbf{Singularitäts-Schutz:} Vermeidet Konfigurationen, bei denen der Roboter unendliche Geschwindigkeiten benötigen würde.
        \item \textbf{Dynamische Anpassung:} Passt die Druckparameter in Echtzeit an, um Materialfluss und Druckqualität zu optimieren.
    \end{itemize}
\end{frame}

\begin{frame}{Kronos Entwicklungsumgebung} % Docker-Container mit ROS, MoveIt, Gazebo
    \begin{alertblock}{Entwicklungsumgebung}
        Kronos wird in einem Docker-Container entwickelt, der ROS (Robot Operating System), MoveIt und Gazebo enthält. 
        Diese Tools ermöglichen die Simulation und Steuerung von Roboterarmen, die Entwicklung von Pfadplanungsalgorithmen und die Integration von Sensoren für die Echtzeit-Überwachung.
    \end{alertblock}

    \vspace{0.5cm}
    \textbf{Vorteile der Docker-basierten Entwicklung:}
    \begin{itemize}
        \item Konsistente Entwicklungsumgebung über verschiedene Systeme hinweg.
        \item Einfache Verwaltung von Abhängigkeiten und Bibliotheken.
        \item Möglichkeit zur schnellen Iteration und Fehlerbehebung durch isolierte Container.
    \end{itemize}
\end{frame}
